% !TeX root = ../thuthesis-example.tex

\chapter{绪论}
\section{国内研究背景}

% injected sources: 01-FA/templates/data/chap01.tex; 01-FA/templates/data/chap02.tex
随着生成式人工智能、学习分析和教育数据挖掘逐步进入高校教学场景，人工智能课程大作业已经不再适合停留在单一算法演示层面。本项目选择 LingoBridge 国际中文学习平台作为工程背景，围绕学生听读练习、录音提交、教师批改、班级学习轨迹沉淀等真实教学流程，构建“学习行为分析与个性化练习推荐”的课程大作业主题。

从国内教育数字化趋势看，中文学习平台、在线课程资源和智慧教学工具正在共同推动语言教学从经验管理转向数据驱动管理。LingoBridge 已具备教师端、学生端和管理端的 MVP 基础能力，能够形成课程、作业、录音、反馈和课堂活动等多类可观测数据。这些数据虽然规模不大，但足以支撑课程大作业所要求的数据预处理、特征构建、模型设计、结果分析和交互展示。

\section{国外研究背景}

% injected sources: 01-FA/templates/data/chap02.tex
在国际语言学习产品中，学习行为追踪、闯关式练习、学习提醒和个性化推荐已经成为常见能力。通用语言学习平台更强调大规模用户增长与订阅转化，而学校内部的国际中文教学更强调班级管理、教师反馈和弱网可用性。两类场景的差异说明，本项目不应简单复制通用 App 的推荐逻辑，而应围绕课堂作业完成率、练习频率、错误类型和学习持续性进行轻量化建模。

因此，本报告将国外成熟语言学习平台的产品思路作为参照，把 LingoBridge 定位为面向教学闭环的轻量 AI 辅助系统。推荐模块不追求复杂黑箱模型，而优先采用可解释的规则特征、统计指标和可演示的机器学习基线，让教师能够看懂推荐依据，学生能够感知练习建议。

\section{研究目的}

% injected sources: prds/json/sprint09-prd-260601-v2.0.json
本大作业的目标是设计并实现一个“基于机器学习的国际中文学习行为分析与个性化练习推荐系统”。系统以 LingoBridge MVP 为基础，在不破坏既有教师端、学生端和管理端主流程的前提下，补充数据分析与智能推荐能力。

具体而言，学生端应展示个人学习轨迹、练习建议和班级排名反馈；教师端应展示学生总览、风险学生识别和作业质量趋势；管理端应作为数据中台，展示全局课程、用户、作业和运行状态指标。通过上述设计，项目能够覆盖课程大作业要求中的数据收集、数据清洗、算法设计、交互设计、编码实现、运行测试和报告撰写等环节。
